% ============================================================================
% Appendix B — Glossary of Arabic Poetry Terms
% ============================================================================
\chapter{Glossary of Arabic Poetry Terms}
\label{app:glossary}

\begin{description}[style=nextline, leftmargin=3cm]

  \item[\ar{البحر} (al-Bahr)] Literally ``the sea''; the term for a classical Arabic meter. Plural: \ar{بحور}.

  \item[\ar{البيت} (al-Bayt)] A verse (literally ``the house''): a complete unit of classical Arabic poetry consisting of two hemistichs (sadr + ajuz).

  \item[\ar{الشطر} (al-Shatr)] A hemistich: one half of a bayt. Each bayt has two shatrs.

  \item[\ar{الصدر} (al-Sadr)] The first hemistich of a bayt (literally ``the chest'').

  \item[\ar{العجز} (al-'Ajuz)] The second hemistich of a bayt (literally ``the rear/end'').

  \item[\ar{التفعيلة} (al-Tafeela)] A metrical foot in Arabic prosody. Plural: \ar{تفعيلات}.

  \item[\ar{الزحاف} (al-Zehaf)] A licensed deviation from the canonical foot pattern in the middle (hashw) of a hemistich. Zehafs are ``permissible'' — they may or may not appear and still yield a valid verse.

  \item[\ar{العلة} (al-'Ella)] A mandatory modification of the aruz or dharb (the last foot of a hemistich). Unlike zehafs, ellas are \textit{fixed} once a particular aruz/dharb is chosen — they must appear consistently.

  \item[\ar{الحشو} (al-Hashw)] The middle feet of a hemistich — all feet except the last. Zehafs apply to hashw.

  \item[\ar{العروض} (al-'Arud)] The last foot of the first hemistich (sadr). Its modification is an \ar{علة عروضية}.

  \item[\ar{الضرب} (al-Dharb)] The last foot of the second hemistich (ajuz). Its modification is an \ar{علة ضربية}.

  \item[\ar{القافية} (al-Qafiya)] The rhyme system — specifically, the cluster of letters at the end of each verse defining the poem's rhyme. Not merely the last sound but a structural unit involving up to five defined letters.

  \item[\ar{الروي} (al-Rawiy)] The primary rhyme consonant — the defining letter of a poem's rhyme.

  \item[\ar{الردف} (al-Radf)] A long vowel (ا، و، ي) immediately preceding the rawiy.

  \item[\ar{التأسيس} (al-Ta'siis)] An alif two positions before the rawiy (with a consonant between).

  \item[\ar{القصيدة} (al-Qasida)] A formal Arabic ode; typically 7+ verses on a consistent meter and rhyme.

  \item[\ar{المطلع} (al-Matla')] The opening verse of a qasida. In classical poetry, the matla' is \ar{مصرَّع} — both hemistichs rhyme with the poem's rawiy.

  \item[\ar{المصراع} (al-Misra')] An individual hemistich.

  \item[\ar{التصريع} (al-Tasri')] The classical convention where the first bayt has both hemistichs rhyming with the rawiy.

  \item[\ar{الغرض} (al-Gharad)] The purpose or thematic genre of a poem (madih, hija', ghazal, etc.). Plural: \ar{أغراض}.

  \item[\ar{النسيب} (al-Nasib)] The amatory prelude of a classical qasida; often a lament at a deserted campsite (\ar{الوقوف على الأطلال}).

  \item[\ar{التخلص} (al-Takhallus)] The transitional section of a qasida, moving from the opening to the main purpose.

  \item[\ar{المديح} (al-Madih)] Panegyric; praise poetry.

  \item[\ar{الهجاء} (al-Hija')] Satire; invective poetry.

  \item[\ar{الغزل} (al-Ghazal)] Love poetry. Two sub-types: \ar{عذري} (Platonic) and \ar{حسي} (sensual).

  \item[\ar{الرثاء} (al-Ritha')] Elegy; mourning poetry.

  \item[\ar{الفخر} (al-Fakhr)] Boasting poetry; self-glorification.

  \item[\ar{الحكمة} (al-Hikma)] Wisdom poetry; philosophical aphorisms.

  \item[\ar{الزهد} (al-Zuhd)] Ascetic poetry; renunciation of worldly pleasures.

  \item[\ar{الخمريات} (al-Khamriyyat)] Wine poetry.

  \item[\ar{الحماسة} (al-Hamasa)] Martial/heroic poetry.

  \item[\ar{المعارضة} (al-Mu'arada)] A response poem using the same meter and rhyme as an existing poem.

  \item[\ar{المساجلة} (al-Musajala)] Poetic competition between two poets exchanging verses on the same meter and rhyme.

  \item[\ar{التضمين} (al-Tadmin)] Incorporating another poet's verse into one's own poem.

  \item[\ar{التشطير} (al-Tashtir)] Inserting a new hemistich before each hemistich of an existing poem.

  \item[\ar{التخميس} (al-Takhmiis)] Adding three new hemistichs to each verse of an existing poem.

  \item[\ar{المحاورة} (al-Muhawara)] Improvised real-time poetic dueling (Hijazi term).

  \item[\ar{القلطة} (al-Qalta)] Same as Muhawara (Najdi/Gulf term).

  \item[\ar{الطريق} (al-Tariq)] The melody-meter path used in Muhawara/Qalta. Plural: \ar{طرق}.

  \item[\ar{العرضة} (al-'Ardah)] A tribal martial performance with poetry, song, and choreography.

  \item[\ar{السامري} (al-Samri)] A Najdi nocturnal performance tradition using Nabati poetry.

  \item[\ar{الزجل} (al-Zajal)] Levantine/Egyptian colloquial strophic poetry tradition.

  \item[\ar{المسحوب} (al-Mashhub)] The dominant Nabati meter; characterised by a dual-rhyme system.

  \item[\ar{الهلالي} (al-Hilali)] The oldest Nabati meter; named for Banu Hilal tribe.

  \item[\ar{القافية المزدوجة} (al-Qafiya al-Muzdawaja)] The dual-rhyme system: distinct rhymes for sadr and ajuz hemistichs in Nabati poetry.

  \item[\ar{شعر التفعيلة} (Shi'r al-Tafila)] Metrical free verse: retains classical foot patterns but uses variable line lengths; no fixed hemistich structure.

  \item[\ar{قصيدة النثر} (Qasidat al-Nathr)] Prose poem: abandons both meter and rhyme; evaluated on imagery density and internal rhythm.

  \item[\ar{التشكيل / التشكيل} (al-Tashkeel)] Full diacritics (harakaat) on Arabic text, required for prosodic scanning.

  \item[\ar{الدائرة العروضية} (al-Da'ira al-'Arudiyya)] One of al-Khalil's five abstract prosodic circles organising related meters.

\end{description}
